\begin{lemma}[A common in-measure limit for the Riesz images of $L^1$ approximating sequences]
Let $f \in L^1(\gamma_n;\mathbb{R})$.  Then there exist an element $g$ of
$\mathcal{M}(\gamma_n;\mathbb{R}^n)$, the space of $\gamma_n$-almost-everywhere
equivalence classes of measurable $\mathbb{R}^n$-valued functions, and a sequence
$(v_k)_{k \in \mathbb{N}}$ in $L^2(\gamma_n)$, such that
\begin{enumerate}
\item $\|\iota v_k\|_{L^1(\gamma_n)} \leq \|f\|_{L^1(\gamma_n)}$ for every $k$, where
$\iota$ is the inclusion of \noderef{ToL1};
\item $\iota v_k \to f$ in $L^1(\gamma_n)$;
\item $R_\gamma v_k \to g$ pointwise $\gamma_n$-a.e., with $R_\gamma$ as in
\noderef{RieszTransform}; and
\item for \emph{every} sequence $(w_k)$ in $L^2(\gamma_n)$ with $\iota w_k \to f$ in
$L^1(\gamma_n)$, the sequence $R_\gamma w_k$ converges to $g$ in measure, i.e.
$\gamma_n\bigl(\{x : \varepsilon \leq |R_\gamma w_k(x) - g(x)|\}\bigr) \to 0$ as
$k \to \infty$, for every $\varepsilon > 0$.
\end{enumerate}
\end{lemma}
\begin{proof}
\emph{Part A: $R_\gamma$ is linear on $L^2(\gamma_n)$.}
By \noderef{RieszTransform}, $R_\gamma v = D L^{-1/2} v$.  Part (1) of
\noderef{FracPowLinear} with $s = -1/2 \leq 0$ says that
$\operatorname{Dom}(L^{-1/2})$ is all of $L^2(\gamma_n)$, so parts (2) and (3) of
\noderef{FracPowLinear} apply to every pair of elements and give
$L^{-1/2}(v+w) = L^{-1/2}v + L^{-1/2}w$ and
$L^{-1/2}(cv) = c\,L^{-1/2}v$ for all $v,w \in L^2(\gamma_n)$ and $c \in \mathbb{R}$.
By \noderef{PseudoInverseSqrt}, $L^{-1/2}v \in \operatorname{Dom}(L^{1/2})$ for every
$v$, and hence $L^{-1/2}v \in W^{1,2}(\gamma_n)$ by \noderef{SqrtDomainSobolev}.
Parts (1) and (2) of \noderef{GaussDLinear} then give
$D(L^{-1/2}v + L^{-1/2}w) = D L^{-1/2}v + D L^{-1/2}w$ and
$D(c\,L^{-1/2}v) = c\,D L^{-1/2}v$.  Composing,
\begin{equation*}
  R_\gamma(v+w) = R_\gamma v + R_\gamma w, \qquad R_\gamma(cv) = c\,R_\gamma v ,
  \tag{A}
\end{equation*}
and in particular, taking $c = -1$, $R_\gamma(v-w) = R_\gamma v - R_\gamma w$.
Also $\iota(v-w) = \iota v - \iota w$, since by \noderef{ToL1} both sides have the
$\gamma_n$-a.e.\ representative $v - w$.

\emph{Part B: the Cauchy-in-measure estimate.}
Let $v,w \in L^2(\gamma_n)$ and $\varepsilon > 0$.  Apply \noderef{MainBound} to the
$L^2$ function $v-w$ at the level $\varepsilon/2 > 0$:
\[
  \tfrac{\varepsilon}{2}\,
    \gamma_n\bigl(\{x : \tfrac{\varepsilon}{2} < |R_\gamma(v-w)(x)|\}\bigr)
    \leq 2\,\|\iota(v-w)\|_{L^1(\gamma_n)} .
\]
By Part A, $R_\gamma(v-w) = R_\gamma v - R_\gamma w$ and
$\iota(v-w) = \iota v - \iota w$, so $\gamma_n$-a.e.\ the function
$R_\gamma(v-w)$ agrees with $R_\gamma v - R_\gamma w$.  Since
$\varepsilon/2 < \varepsilon$, the set $\{x : \varepsilon \leq |R_\gamma v(x) -
R_\gamma w(x)|\}$ is therefore contained, up to a $\gamma_n$-null set, in
$\{x : \varepsilon/2 < |R_\gamma(v-w)(x)|\}$, and monotonicity of $\gamma_n$ gives
\[
  \tfrac{\varepsilon}{2}\,
  \gamma_n\bigl(\{x : \varepsilon \leq |R_\gamma v(x) - R_\gamma w(x)|\}\bigr)
    \leq 2\,\|\iota v - \iota w\|_{L^1(\gamma_n)} .
\]
Dividing by $\varepsilon/2 > 0$,
\begin{equation*}
  \gamma_n\bigl(\{x : \varepsilon \leq |R_\gamma v(x) - R_\gamma w(x)|\}\bigr)
    \leq \frac{4\,\|\iota v - \iota w\|_{L^1(\gamma_n)}}{\varepsilon} .
  \tag{B}
\end{equation*}

\emph{Part C: the truncations.}
Write $F$ for a fixed representative of $f$ and put
$\tau_k(x) = \max\{-k, \min\{F(x), k\}\}$ for $k \in \mathbb{N}$.  Each $\tau_k$ is
measurable, being built from $F$ by two lattice operations with constants.  For every
$x$ and $k$ we have, since $k \geq 0$,
\begin{equation*}
  |\tau_k(x)| \leq k \qquad\text{and}\qquad |\tau_k(x)| \leq |F(x)| .
  \tag{C1}
\end{equation*}
For the first: $-k \leq \max\{-k,\min\{F(x),k\}\} \leq \max\{-k,k\} = k$.  For the
second: $-|F(x)| \leq \min\{F(x),k\}$ because $-|F(x)| \leq F(x)$ and
$-|F(x)| \leq 0 \leq k$, whence $-|F(x)| \leq \tau_k(x)$; and
$\tau_k(x) \leq |F(x)|$ because $-k \leq 0 \leq |F(x)|$ and
$\min\{F(x),k\} \leq F(x) \leq |F(x)|$.
Since $\gamma_n$ is a probability measure (\noderef{GaussIsProb}), the bound
$|\tau_k| \leq k$ makes $\tau_k$ an element of $L^2(\gamma_n)$; call it $\tilde v_k$.
Moreover, for every $x$, as soon as $k \geq |F(x)|$ we have
$-k \leq F(x) \leq k$, hence $\tau_k(x) = F(x)$; so
$\tau_k(x) \to F(x)$ for every $x$, the sequence being eventually constant.
The second bound in (C1) gives, by monotonicity of the integral,
$\|\iota \tilde v_k\|_{L^1(\gamma_n)} = \int|\tau_k|\,d\gamma_n
  \leq \int |F|\,d\gamma_n = \|f\|_{L^1(\gamma_n)}$,
which is conclusion (1) for this sequence.  Finally
$|\tau_k(x) - F(x)| \leq |\tau_k(x)| + |F(x)| \leq 2|F(x)|$ by (C1), the dominating
function $2|F|$ is integrable because $f \in L^1(\gamma_n)$, and
$|\tau_k(x) - F(x)| \to 0$ for every $x$; so dominated convergence gives
$\|\iota\tilde v_k - f\|_{L^1(\gamma_n)} = \int|\tau_k - F|\,d\gamma_n \to 0$, i.e.
$\iota\tilde v_k \to f$ in $L^1(\gamma_n)$, which is conclusion (2).

\emph{Part D: a rapidly converging subsequence.}
Since $\|\iota\tilde v_k - f\|_{L^1(\gamma_n)} \to 0$, for each $j \in \mathbb{N}$ we
may choose an index $\sigma(j)$ with
$\|\iota\tilde v_{\sigma(j)} - f\|_{L^1(\gamma_n)} \leq \tfrac12 8^{-j}$.  Put
$v_j = \tilde v_{\sigma(j)}$.  Conclusion (1) is inherited, and
$\|\iota v_j - f\|_{L^1(\gamma_n)} \leq \tfrac12 8^{-j} \to 0$ gives conclusion (2)
for $(v_j)$.  By the triangle inequality and $8^{-(j+1)} \leq 8^{-j}$,
\begin{equation*}
  \|\iota v_j - \iota v_{j+1}\|_{L^1(\gamma_n)}
    \leq \|\iota v_j - f\|_{L^1(\gamma_n)} + \|\iota v_{j+1} - f\|_{L^1(\gamma_n)}
    \leq \tfrac12 8^{-j} + \tfrac12 8^{-(j+1)} \leq 8^{-j} .
  \tag{D}
\end{equation*}

\emph{Part E: the almost-everywhere limit.}
Put $E_j = \{x : 2^{-j} \leq |R_\gamma v_j(x) - R_\gamma v_{j+1}(x)|\}$.  By (B) with
$\varepsilon = 2^{-j}$ and then (D),
\[
  \gamma_n(E_j) \leq \frac{4\,\|\iota v_j - \iota v_{j+1}\|_{L^1(\gamma_n)}}{2^{-j}}
    \leq 4\cdot 2^{j}\cdot 8^{-j} = 4\cdot 4^{-j} .
\]
The series $\sum_j 4\cdot4^{-j}$ converges, so $\sum_j \gamma_n(E_j) < \infty$.  By the
Borel--Cantelli lemma, for $\gamma_n$-a.e.\ $x$ the index set $\{j : x \in E_j\}$ is
finite.  Fix such an $x$ and let $B$ be an upper bound for that finite set.  For every
$j \geq B+1$ we then have $x \notin E_j$, i.e.
$|R_\gamma v_j(x) - R_\gamma v_{j+1}(x)| < 2^{-j}$.  Hence the series
$\sum_j |R_\gamma v_j(x) - R_\gamma v_{j+1}(x)|$ has all but finitely many terms
dominated by the convergent geometric series $\sum_j 2^{-j}$, so it converges; a
sequence in the complete space $\mathbb{R}^n$ whose consecutive distances form a
convergent series is Cauchy, hence converges.  Define $g(x)$ to be this limit at every
$x$ for which $(R_\gamma v_j(x))_j$ converges, and $g(x) = 0$ elsewhere.  By
construction $R_\gamma v_j \to g$ pointwise $\gamma_n$-a.e., which is conclusion (3);
and $g$, being a $\gamma_n$-a.e.\ pointwise limit of the measurable functions
$R_\gamma v_j$, is measurable, so it determines an element of
$\mathcal{M}(\gamma_n;\mathbb{R}^n)$.

\emph{Part F: every approximating sequence has the same limit in measure.}
Since $\gamma_n$ is a finite measure (\noderef{GaussIsProb}), the a.e.\ convergence of
conclusion (3) implies $R_\gamma v_k \to g$ in measure; that is, for every
$\delta > 0$,
\begin{equation*}
  \gamma_n\bigl(\{x : \delta \leq |R_\gamma v_k(x) - g(x)|\}\bigr) \longrightarrow 0 .
  \tag{F}
\end{equation*}
Now let $(w_k)$ be any sequence in $L^2(\gamma_n)$ with $\iota w_k \to f$ in
$L^1(\gamma_n)$, and fix $\varepsilon > 0$.  By the triangle inequality in
$\mathbb{R}^n$, if $\varepsilon \leq |R_\gamma w_k(x) - g(x)|$ then at least one of
$|R_\gamma w_k(x) - R_\gamma v_k(x)|$ and $|R_\gamma v_k(x) - g(x)|$ is at least
$\varepsilon/2$; hence
\[
  \{x : \varepsilon \leq |R_\gamma w_k(x) - g(x)|\}
    \subseteq \{x : \tfrac{\varepsilon}{2} \leq |R_\gamma w_k(x) - R_\gamma v_k(x)|\}
      \cup \{x : \tfrac{\varepsilon}{2} \leq |R_\gamma v_k(x) - g(x)|\} ,
\]
and by monotonicity and subadditivity of $\gamma_n$,
\[
  \gamma_n\bigl(\{x : \varepsilon \leq |R_\gamma w_k(x) - g(x)|\}\bigr)
    \leq \gamma_n\bigl(\{x : \tfrac{\varepsilon}{2}
          \leq |R_\gamma w_k(x) - R_\gamma v_k(x)|\}\bigr)
      + \gamma_n\bigl(\{x : \tfrac{\varepsilon}{2}
          \leq |R_\gamma v_k(x) - g(x)|\}\bigr) .
\]
The second summand tends to $0$ by (F) with $\delta = \varepsilon/2$.  For the first,
$\|\iota w_k - \iota v_k\|_{L^1(\gamma_n)} \to 0$, because both $\iota w_k$ and
$\iota v_k$ converge to $f$ in $L^1(\gamma_n)$; so (B) with $\varepsilon/2$ bounds the
first summand by
$8\,\|\iota w_k - \iota v_k\|_{L^1(\gamma_n)}/\varepsilon \to 0$.  Both summands are
nonnegative, so their sum tends to $0$ and the left-hand side, squeezed between $0$ and
that sum, tends to $0$ as well.  This is conclusion (4).
\end{proof}
